\subsection{Arithmetic at the terminal set of auxiliary primes}
\label{sec:terminal-selmer}
We give the length estimate used in \S\ref{sec:selmer-polynomial-descent}
for a Selmer datum satisfying the following hypotheses. Write $T=O^n$,
$G=G_E$ for a CM field $E$, and $S$ for a finite set containing the primes over $\ell$, outside which $\rho$ is unramified. Fix the torsion level $M$ and take $N>M$. Fix the common element $g_0$ in inertia at a bad place away from $\ell$ and a common open neighborhood imposing the prescribed congruences there. At each $v\mid\ell$
specify a finite free lattice $Z_v$ of local integral cocycles containing
the coboundaries. Impose on the torsion cocycles all the homogeneous
linear relations among evaluations valid on $Z_v$. Require
\[
 \sum_{v\mid\ell}\operatorname{rank} Z_v
 \leq (\#\{v\mid\ell\}+r_2(E))n .
\]
Assume $\rho(G)$ spans $M_n(O)$ over $O$. Suppose there is
$\sigma$ in inertia at a prime of $S$ away from $\ell$ such that
$\rho(\sigma)=\rho(\sigma^{-c})=\xi I$, with $\xi^2-1$ a unit.
Add the normalization $z(\sigma)=0$ to the global cocycles.
For $W=T/\pi^j$, let $W^*=\operatorname{Hom}(W,\mu_{\ell^\infty})$
be the Tate dual with the contragredient action and let $W^{*c}$ carry
the action pulled back by $g\mapsto g^c=cgc^{-1}$. In additive
notation for characters, scalars in $O$ act by precomposition.
Define $G_T=\sup_j\lambda_O\operatorname{Hom}_G(W,W^{*c})$.
For $v\nmid\ell$ let $b_v$ be the supremum over $j$ of the positive
part of $\lambda_O C_v(W)-\lambda_O W$, where $C_v(W)$ consists of
restrictions to $G_{E_v}$ of all continuous global cocycles with values
in $W$ and unrestricted ramification. Put $L_T=\sum_{v\in S,\ v\nmid\ell} b_v$.
Suprema over all nonzero scalar ideals give the same definitions.
The argument presumes finiteness of these terms as in the datum.

\paragraph{Producing allowable primes.}
Put $q=\ell^N$. The following construction also explains the choice of
the first prime $w_0$. Fix finitely many conjugation-stable places $B$
containing the places over $\ell$, at which the new character is to be
locally trivial. Let $D$ be a finite set disjoint from $B$ and let
$g$ belong to the normal subgroup generated by inertia at $D$. For any open
neighborhood of $g$, one can choose a degree-one prime $w\ne\bar w$ off a
prescribed finite exceptional set, a Frobenius lift in that neighborhood
and an anticyclotomic character $h:G\to\mathbb Z/q$ trivial locally at $B$,
unramified off $w,\bar w$, and onto on inertia at $w$. Its tame generator
may be chosen to have image $1$ and its Frobenius lift to have image $0$.

Indeed enlarge $B$ to a finite set $B'$ disjoint from $D$ and generating
the ideal class group: choose representatives prime to the given excluded
set and take their prime divisors. The extension of $E$ formed by adjoining $\mu_q$
and $q$th roots of generators of the finitely generated $B'$-unit group
is finite Galois and unramified at $D$, since these generators are units
there and $q$ is invertible. Its Galois kernel contains $g$. Apply
Chebotarev~\cite[Theorem~8.31]{milneANT} to the compositum with a finite Galois quotient specifying
the neighborhood: one can choose a place above $w$ with the prescribed
Frobenius element in that quotient. The set of possible primes has
positive Dirichlet density, hence contains degree-one primes since
the primes of higher residue degree have convergent norm-reciprocal sum.
Avoid primes ramified over $\mathbb Q$. A degree-one unramified prime
cannot be fixed by $c$, since its decomposition group in $E/E^+$ is trivial.
We also exclude $B'$ and any other specified finite set.
The residue field at $w$ contains $\mu_q$ and the $B'$-units reduce
to $q$th powers. A surjective discrete logarithm from its multiplicative
group to $\mathbb Z/q$ therefore kills every $B'$-unit.
On finite ideles which are units off $B'$ take the logarithm of
the residue at $w$. This map descends to a character of all finite idele
classes: principal ideles and this subgroup generate, and their
intersection consists of principal $B'$-units. Its kernel is open.
Precompose with $a\mapsto a/\bar a$ and use global Artin reciprocity.
The result has the claimed ramification and is trivial on the full local
multiplicative groups at $B$ because $B$ is conjugation-stable and
$w,\bar w\notin B$. It surjects on inertia since $w\ne\bar w$.
Multiply the chosen Frobenius lift by a power of its tame generator to
give character value zero. This preserves its image in the prescribed
finite quotient since $w$ can be chosen unramified there.

Start with $B=\{v\mid\ell\}$, and use $g_0$ in the prescribed bad
inertia away from $\ell$. Take $w_0$ with Frobenius close enough to
$g_0$ for all the finitely many congruences on the data and the
Stickelberger systems. Exclude the rational primes below the whole prescribed bad set when choosing $w_0$, and exclude the rational prime below $w_0$ in every subsequent choice. Take all slopes zero there. For extensions of
$J$, put $B=\{v\mid\ell\}\cup(J\setminus\{w_0\})\cup
\overline{(J\setminus\{w_0\})}$. For the datum being shortened let
$H$ be the subgroup generated by $[g,\sigma]$ and $[g,\sigma^c]$,
$g\in G$, using $[a,b]=aba^{-1}b^{-1}$. It is normal, since
$a[g,b]a^{-1}=[ag,b][a,b]^{-1}$, and stable under $c$.
Its image under $\rho$ and every abelian character is trivial. For
any $h\in H$, take as target $g=\gamma_0 h$, where $\gamma_0$ is the
tame generator at $w_0$. This lies in the normal subgroup generated by
inertia at $w_0$ and the primes for $\sigma,\sigma^c$, disjoint from $B$.
At this target the old characters take values $1$ for $h_{w_0}$ and $0$
for the others. Both $g,g^c$ act trivially modulo $\ell^M$ and on
$\mu_q$. These are open conditions and so can be imposed on the new
Frobenius. The new character is trivial at the old later primes, even
on their Frobenius elements, by local triviality at $B$. We can keep
the new prime and its conjugate away from $S_0$, all old primes and
the rational prime below $w_0$. This extends all the local specifications
of \S\ref{sec:selmer-polynomial-descent}, for arbitrary slope on the
chosen datum. In particular, allowable Frobenii approximate every
$\gamma_0 h$ with $h\in H$ arbitrarily closely in $G$. The old characters are held constant and the actions on $T/\ell^M$, its conjugate and $\mu_q$ held trivial by shrinking the neighborhood.

\paragraph{Comparing a Selmer group with its dual.}
Fix $j\le e$. Define local subgroups $\mathcal L_v\subset H^1(E_v,W)$
as follows. At $v\mid\ell$ use cocycles satisfying the relations of
$Z_v$, modulo coboundaries. At $v\in S$ away from $\ell$ use the image
of $C_v(W)$. At a later auxiliary prime and its conjugate use the
cocycles satisfying the tie, modulo coboundaries. Elsewhere use
unramified classes, including at $w_0,\bar w_0$. The local conditions
contain coboundaries before quotienting: at a later prime inertia acts
trivially, and a nonzero slope requires $\rho(\phi)=1$ on $W$.
Every global class has a unique representative vanishing at $\sigma$,
because $\xi-1$ is invertible. Consequently the normalized tied cocycles
are exactly $H^1_{\mathcal L}(E,W)$. Use the orthogonal complements
$\mathcal L_v^\perp$ for the local Tate pairings to define its dual
$Y=H^1_{\mathcal L^\perp}(E,W^*)$. The finite-module Greenberg--Wiles
formula (A. Smith, \emph{Galois cohomology and the Greenberg--Wiles
formula}, Theorem 1.8~\cite{smith}) gives
\[
 \lambda_O Z_{i,j}(J)-\lambda_O Y
 = \sum_{v\text{ finite}}
       (\lambda_O\mathcal L_v-\lambda_O H^0(E_v,W))
       -r_2(E)\lambda_O W .
\]
Here $H^0(E,W)$ and $H^0(E,W^*)$ vanish since $\sigma$ acts by $\xi$
and $\xi^{-1}$ respectively, the cyclotomic character being trivial at
$\sigma$. Taking logarithms of the cardinality formula to base $|k|$
gives these lengths. At an unramified place the local difference is
zero. At any other finite place, if $D_v$ is the specified local
cocycle group, the difference equals $\lambda_O D_v-\lambda_O W$,
by the exact sequence defining coboundaries. At a tied place a cocycle
is determined by its values on inertia and Frobenius, and the former
value is determined by the latter. Thus $\lambda_O D_v\leq\lambda_O W$.
At a place in $S$ away from $\ell$, the excess costs at most $b_v$.

For $v\mid\ell$ let $r_v=\operatorname{rank} Z_v$.
All scalar coordinate evaluations on $Z_v$ generate a free
submodule of $Z_v^\vee$ of rank at most $r_v$. Each local solution to
the defining linear relations gives a homomorphism from this submodule to
$O/\pi^j$, by sending a coordinate evaluation to the corresponding
coordinate of the solution. Every dependence on the evaluations is a
relation by definition. This map on solutions is injective, so
$\lambda_O D_v\le jr_v$. The stated bound on the sum of ranks offsets the
possible archimedean loss. We obtain
\[
                    \lambda_O Z_{i,j}(J)\leq\lambda_O Y+L_T .
\]
The choice of $C_v(W)$ at old bad places matters: using all local
cocycles there could enlarge the local error in the definition of $L_T$.

\paragraph{Controlling the dual at a terminal set.}
Identify the underlying constant group of roots of unity of order dividing
$\ell^M$ with the corresponding subgroup of $\mathbb Q/\mathbb Z$
using a compatible primitive root; all Galois actions retain their
cyclotomic factors. Suppose $z\in Z_{i,j}(J)$ has some value with a unit
coordinate modulo $\pi^j$. On $H$, $z$ and each dual cocycle $y$
are homomorphisms. The span of $z(H)$ is $W$: it is stable under
$\rho(G)$, and $z([g,\sigma])=(1-\xi)z(g)$ since $z(\sigma)=0$.
A vector with a unit coordinate generates $W$ under the full matrix
algebra. The same spanning assertion holds for $z(H^c)$.

For $t\in O$ define the biadditive function
\[
 B_t(h,h')= y(h')(t z(h))-y(h^c)(t z(h'^c)),\qquad h,h'\in H .
\]
We show $B_t(h,h)=0$. The restrictions of $y$ to inertia at $w_0$ and
$\bar w_0$ are zero: the dual of the unramified condition there is
unramified, and inertia acts trivially. The same holds for $z$ by the
zero tie. Their values at $\gamma_0 h$ and $(\gamma_0h)^c$ therefore
coincide with those at $h,h^c$. Choose a new auxiliary prime as in the
preceding construction, with Frobenius $\phi$ sufficiently close to
$\gamma_0h$ that these values of $y,z$ and their conjugates agree.
They are locally constant. Both $\phi$ and $\phi^c$ act trivially
on $W,W^*$. The terminal inclusion produces a cocycle $z'$ allowed to
ramify at this new pair with
\[
 z'(\gamma)=z(\phi),\qquad z'(\gamma^{-c})=z(\phi^{-c})=-z(\phi^c).
\]
The sum of the local invariants of the global cup product of $tz'$ and
$y$ vanishes. At old places the classes are orthogonal by construction;
away from the new pair and the old places both are unramified. At the
new pair the dual class is unramified. For trivial action the local
pairing of a cocycle and an unramified dual class evaluates the latter
on Frobenius against the value of the former on normalized tame inertia
(up to a common sign from the cup-product order). In detail, cup
product followed by the tame residue contracts the restriction to
$\operatorname{Hom}(I^{\rm tame},W)$ with the unramified Frobenius
value; the normalization sends the tame generator acting on a
$q$th root of a uniformizer by the chosen primitive $q$th root to
$1$. For a cyclic summand of order $m\mid\ell^M$, identify it locally with $\mu_m$ using the compatible root. A Kummer class of $a\in E_w^\times$ has tame value $v_w(a)\bmod m$ in this normalization: the Kummer extension of a unit is unramified since $m$ is prime to the residue characteristic. Pairing with an unramified character evaluates that character on the local Artin symbol of $a$, hence on $v_w(a)$ times Frobenius. Additivity over the cyclic summands gives the assertion for $W$. At the conjugate place the normalized generator is the inverse
conjugate because $c$ inverts roots of unity. Replacing the two normalized
generators by $\gamma,\gamma^{-c}$ multiplies both evaluations by the
same unit modulo $\ell^M$. Thus the vanishing reads
$y(\phi)(t z(\phi))-y(\phi^c)(t z(\phi^c))=0$, as claimed.

Polarizing gives $B_t(h,h')+B_t(h',h)=0$. Conjugate its first argument
by $\sigma$. The two transformation rules are
\[
 B_t(\sigma h\sigma^{-1},h')=B_{\xi t}(h,h'),\qquad
 B_t(h',\sigma h\sigma^{-1})=B_{\xi^{-1} t}(h',h).
\]
Indeed $\sigma$ and $\sigma^c$ have opposite scalar actions on $W$,
and the dual actions invert these scalars. Substituting $\xi t$ in
the polarized identity with conjugated argument and subtracting the
original identity gives $B_{(\xi^2-1)t}(h,h')=0$. Thus $B_t=0$ for all $t$.

The assignment $z(h)\mapsto y(h^c)$ extends uniquely to an
$O$-linear map $\eta_y:W\to W^{*c}$. To see well-definedness, take a
relation $\sum_j a_j z(h_j)=0$. Evaluating
$\sum_j a_j y(h_j^c)$ on $t z(h'^c)$ gives zero for all $h',t$, by
$B_{a_jt}(h_j,h')=0$. These vectors span $W$ additively, including
all their scalar multiples. Conjugating $h$ by $g\in G$ gives
$\eta_y(\rho(g)z(h))=g^c\cdot y(h^c)$ since $H$ acts trivially
on both modules. Hence $\eta_y$ is equivariant. Coboundaries vanish
on $H$, so this is a linear map on dual classes. If $\eta_y=0$,
$y$ vanishes on $H$. Normalize $y(\sigma)=0$ by its unique
coboundary since $\xi^{-1}-1$ is a unit. Then
$y([g,\sigma])=(1-\xi^{-1})y(g)$ shows $y=0$.
Consequently $\lambda_O Y\leq G_T$ whenever such a $z$ exists.

Finally let $j$ be the largest exponent of the cyclic submodules of
the tied module at level $e$ (so these cyclic orders are $\pi^j$).
If $j=0$ the module is zero. Otherwise every value of every such cocycle
is in $\pi^{e-j}W_e$, for $W_e=T/\pi^e$. Identify $T/\pi^j$ with
this submodule by multiplication by $\pi^{e-j}$. Dividing values in
this identification embeds the top-level module in $Z_{i,j}(J)$:
all equations are homogeneous over $O$ and all inertia restrictions
are preserved. A cocycle of maximal order yields a unit coordinate
somewhere. The terminal inclusion was obtained for every level.
Applying the preceding argument at $j$ now gives
$\lambda_O Z_{i,e}(J)\leq G_T+L_T$ for every datum.
This closes the terminal estimate for the indicated Selmer data.
